\section{데이터와 증거 수준}
\label{sec:data}

\subsection{평가 데이터의 출처}
본 논문은 nuScenes 기반 주행 레코드에 추론 주석이 결합된 CoC-Nusc 형식 데이터를 사용한다. 원천 주행 기록은 nuScenes v1.0-trainval 계열이며, 공식 nuScenes는 카메라, 라이다/레이더, 자차 자세 및 선택적 CAN 버스 확장을 제공한다 \cite{caesar2020nuscenes,nuscenes2026website}. 그러나 본 논문의 로컬 CoC-Nusc 실험에서 직접 사용한 물리 증거는 변환된 카메라 패키지, 장면 내부 자차 움직임 parquet, 추론 주석이다. 제3자 변환 과정은 nuScenes 장면을 PhysicalAI-AV 호환 디렉터리와 타임스탬프 관례로 재구성하고, NVIDIA 공개 CoC 자동 레이블러 형식의 행동 주석을 생성한다 \cite{nvidia2026autolabeler}. NVIDIA는 자동 레이블러 소프트웨어를 제공하지만, 본 연구에서 사용한 CoC-Nusc 변환 데이터셋을 생성하거나 인증한 것은 아니다.

따라서 본 논문에서 ``CoC-Nusc''는 공식 nuScenes 원본 자체를 의미하지 않는다. 이는 nuScenes 주행 기록을 기반으로 한 변환 평가 데이터이며, 추론 문장과 품질 점수는 VLM/LLM 기반 자동 레이블링 절차에서 생성된 약한 주석이다. 본 논문의 모델 생성기와 \uppaal 관찰자는 이 변환 데이터를 추가로 해석하여 파싱된 행동 의도, 자차 움직임 응답 증거, 연구자가 선택한 후보 시간 계약, 변이 모델을 도출한다. 표~\ref{tab:provenance-stack}에서는 이 계층을 구분한다.

\begin{table*}[t]
\centering
\caption{평가 데이터의 출처 계층.}
\label{tab:provenance-stack}
\footnotesize
\begin{tabularx}{\textwidth}{L{0.15\textwidth}L{0.18\textwidth}Y L{0.17\textwidth}}
\toprule
Layer & Artifact & Role in this work & Evidence status \\
\midrule
nuScenes/Motional & v1.0-trainval, optional CAN extension & Multi-camera scenes, timestamps, ego pose/speed; actuator commands are not used here & Source driving log \\
CoC-Nusc transform & PhysicalAI-compatible data & Resampling, camera packaging, ego-motion parquet, clip index & Derived compatible data \\
NVIDIA & CoC auto-labeler & Meta-actions, keyframes, VLM prompt, validator & Public labeling software \\
Qwen/Qwen3.5-35B-A3B & Generated reasoning and self-evaluation & One-sentence behavior annotation, quality score & Machine-generated weak label \\
This work & Compiler/observers & Intent parsing, response detection, contracts, mutations & Derived verification evidence \\
\bottomrule
\end{tabularx}
\end{table*}

표~\ref{tab:provenance-stack}의 출처 구분은 이후 실험 결과를 해석할 때 기준이 된다. 예를 들어 CoC 문장과 품질 점수는 원천 센서값이 아니라 기계 생성 약한 레이블이며, 본 연구의 의도 분류와 응답 시작 시각은 다시 한 번 규칙으로 도출된 검증 증거이다. 따라서 어떤 모델이 통과하더라도 그 통과는 원천 주행 기록 전체의 안전성을 의미하지 않고, 해당 계층에서 허용되는 주장으로 제한된다.

\subsection{추론 주석 생성 절차}
CoC 주석은 이미 기록된 영상과 자차 궤적에서 메타 행동을 도출하고, 행동 전환 근처의 키프레임을 선택한 뒤, 영상 클립, 자차 궤적, 메타 행동 힌트를 VLM/LLM에 입력하여 인과적 행동 주석을 생성하는 절차를 따른다 \cite{nvidia2026autolabeler}. 즉 온라인 제어기가 추론을 출력한 것이 아니라, 기록된 미래 움직임을 포함한 클립에서 생성된 회고적 보조 레이블이다.

\begin{figure}[t]
\centering
\begin{tikzpicture}[
  node distance=5mm,
  box/.style={draw, rounded corners=2pt, align=center, minimum height=8mm,
    text width=0.88\columnwidth, font=\footnotesize},
  obs/.style={box, fill=observed!10},
  der/.style={box, fill=derived!10},
  ass/.style={box, fill=assumed!10},
  input/.style={box, fill=black!3},
  flow/.style={-{Latex[length=2mm]}, thick}
]
\node[obs] (a) {Driving video + recorded ego trajectory};
\node[der, below=of a] (b) {Recorded-trajectory-conditioned meta-action extraction};
\node[der, below=of b] (c) {Keyframe selection around action changes};
\node[input, below=of c] (d) {VLM/LLM inputs: recorded video/trajectory + derived meta-action hints};
\node[der, below=of d] (e) {Retrospective CoC reasoning sentence and quality score};
\node[der, below=of e] (f) {This work: parsed intent, audit-input acceptance, ego-motion response evidence};
\draw[flow] (a) -- (b);
\draw[flow] (b) -- (c);
\draw[flow] (c) -- (d);
\draw[flow] (d) -- (e);
\draw[flow] (e) -- (f);
\end{tikzpicture}

\caption{CoC 추론 주석 생성과 본 논문 검증 입력으로의 변환. 자동 레이블링은 기록된 궤적에서 메타 행동을 먼저 도출하므로, 생성 문장은 온라인 인과 증명이 아니라 회고적 추론 지도 신호로 해석된다.}
\label{fig:annotation-flow}
\end{figure}

학습 관점에서 추론 주석은 기존 E2E 궤적 샘플을 확장하는 보조 목표이다. 일반적인 E2E 샘플이 관측 이력 $x_t$와 미래 궤적 $y_t$의 쌍이라면, 추론 보강형 샘플은 인과 문장 $z_t$를 추가해 $(x_t,z_t,y_t)$ 또는 쌍 목표 $(z_t,y_t)$를 구성한다. Alpamayo는 이런 방향의 대표 사례로 CoC 문장과 궤적 예측을 연결해 긴 꼬리 주행 일반화를 목표로 한다 \cite{wang2025alpamayo}. 그러나 본 논문은 실제 E2E 모델을 미세조정하지 않고, 생성된 추론 주석과 기록 자차 움직임이 검증 가능한 에피소드 궤적을 이루는지 다룬다.

그림~\ref{fig:annotation-flow}에서 중요한 지점은 메타 행동 추출이 기록된 궤적에 조건화되어 CoC 생성보다 먼저 수행된다는 것이다. 즉 CoC 문장은 주행 중 온라인으로 산출된 제어 판단이 아니라, 기록된 영상과 궤적을 함께 본 뒤 생성된 회고적 설명이다. 따라서 CoC와 같은 궤적 사이의 일치는 독립적인 인과 증거가 아니다. 이 때문에 본 논문은 CoC 문장을 직접 제어 명령으로 해석하지 않고, 검증기가 확인할 파싱된 추상 행동 의도와 연구자 선택 후보 시간 계약의 입력으로만 사용한다.

\subsection{감사 입력 선별 흐름}
청크~0에서의 감사 입력은 다음의 축소 흐름으로 구성된다.
\begin{quote}
\small
403 CoC 주석 $\rightarrow$ 290 기계 품질 통과 주석 $=134$ 방향 명시 주석 사건별 후보 계약 $+156$ 제외 $\rightarrow$ 130 속도 탐지기 사례 $+$ 4 이미 충족된 사례.
\end{quote}
290개 품질 통과 사건 중 156개는 영어 정규식 기반 의도 파서에서 방향이 모호하거나 비방향적이거나 복합 의도로 판정되어 제외했고, 추론 범위를 방향이 명시된 134개 사건으로 제한했다. 이 134건은 단일 이벤트 선별의 \emph{주석 사건별 후보 계약} 134개이다. 반복 사건끼리 같은 응답 시작을 공유할 수 있으므로 독립 응답 134개를 뜻하지 않으며, \preserve/\reset 상태 분석은 선택한 에피소드 사례에만 적용했다. 즉 $134=125+4+5$이며, 기본 분석 설정에서는 125개 속도 변화 계약 통과, 4개 이미 충족된 계약 통과, 5개 검토 후보로 분기한다. 이 수는 기록 전체 또는 403개 주석 전체에 대한 안전 판정이 아니다.

초기 기본 탐지기 선별에서는 검토 후보가 8건이었다. 이후 STOP/YIELD 사건 시점의 속도가 $0.5$~m/s 이하이면 이미 충족된 것으로 처리하는 규칙을 사후 추가하여 \texttt{scene-0046@10000000}, \texttt{scene-0159@4500000}, \texttt{scene-0159@4900000}을 재분류했고 후보가 5건이 되었다. 이 $0.5$~m/s 임계값은 요구사항에서 독립적으로 도출한 값이 아니라 해당 코호트에 맞춘 탐색적 보정 논리이다.

여기서 \texttt{quality\_pass}는 Qwen/Qwen3.5-35B-A3B가 생성한 \emph{기계 생성 감사 입력 수락 신호}이다. 상류 README에는 \texttt{Qwen/Qwen3.5-35B-A3B}라는 모델명은 기록되어 있으나, 불변 모델 스냅샷·태그·해시와 자동 레이블러 커밋은 보존되어 있지 않다. 따라서 특정 리비전을 추정하지 않으며, 이 누락을 재현성 한계로 취급한다. 이는 인간 정답(human gold)이나 주석 내용의 사실성을 보증하는 표지가 아니다. 다섯 검토 후보에 대한 시각 점검도 각 후보의 전방 카메라 네 프레임을 사용한 AI 시각 분류(triage)이며, 독립적인 인간 검증이 아니다. 이 점검은 검토 우선순위를 돕는 보조 문맥일 뿐 CoC--궤적 일치의 독립 인과 근거나 물리 안전성 결론을 제공하지 않는다.

\subsection{데이터 한 건의 실제 구조}
CoC-Nusc 추론 파일은 \texttt{data/baseline/coc\_nusc/reasoning/ood\_reasoning.parquet}이며, \texttt{clip\_id}별 행 안에 \texttt{events} JSON 문자열이 저장된다. 각 사건이 하나의 추론 보강형 레코드이다. 핵심 필드는 \texttt{clip\_id}, \texttt{event\_start\_timestamp}, \texttt{cot}, \texttt{quality}, \texttt{quality\_pass}이다.

\noindent\begin{minipage}{\columnwidth}
\begin{lstlisting}[caption={scene-0001의 CoC 사건 레코드 예시. 점수는 모델이 생성한 검증 신호이다.},label={lst:coc-record}]
{
  "clip_id": "scene-0001",
  "event_start_timestamp": 2500000,
  "cot": "Decelerate to maintain a safe distance from
          the lead truck ahead while navigating through
          the construction zone.",
  "quality": {
    "visual_evidence_score": 0.90,
    "behavior_match_score": 0.95,
    "decision_supported_by_video": 0.75,
    "causal_consistency_score": 0.70
  },
  "quality_pass": true
}
\end{lstlisting}
\end{minipage}

같은 \texttt{clip\_id} 안의 CoC 주석 사건 목록을 시간순으로 정렬하면 에피소드 궤적이 된다. \texttt{scene-0001}에는 2.5초, 3.5초, 4.8초에 트럭/공사 구간 관련 감속 추론이 반복되고, 6.3초에는 좌측 회피 추론이 나타나며, 9.3초에는 낮은 품질의 사건이 나타난다. 자차 움직임 보조 파일은 \texttt{data/baseline/coc\_nusc/labels/egomotion/scene-0001.egomotion.parquet}이며 타임스탬프, 위치, 쿼터니언, 속도, 가속도, 곡률 등을 포함한다. 본 논문은 $v=\sqrt{v_x^2+v_y^2}$로 자차 속도를 계산하고, CoC 주석 사건 이후의 지속적 속도 변화를 자차 움직임 응답 증거로 도출한다.

\subsection{증거 수준 구분}
모든 사건과 전환은 네 수준으로 구분한다. 이 구분은 ``검증 도구가 실행되었다''는 사실이 ``차량이 안전하다''는 주장으로 확대 해석되는 것을 방지하기 위해 필요하다.

\begin{table}[t]
\centering
\caption{증거 수준과 허용되는 주장.}
\label{tab:evidence-level}
\footnotesize
\begin{tabularx}{\columnwidth}{L{0.24\columnwidth}YY}
\toprule
Level & Example & Allowed claim \\
\midrule
Observed evidence & \texttt{cot}, timestamp, camera frame, ego-motion row & The record exists in the data \\
Derived evidence & Intent, ego speed, response onset, latency & The response is computed under stated rules \\
Synthetic assumption & $D$, repeat semantics, detector setting & Candidate requirements and sensitivity are tested \\
Synthetic mutation & No response, clock reset, overwrite, orphan response & Queries detect injected known faults \\
\bottomrule
\end{tabularx}
\end{table}

표~\ref{tab:evidence-level}에서는 표~\ref{tab:provenance-stack}의 출처 계층을 주장 수준으로 다시 번역한다. 관측 증거는 ``데이터에 그렇게 기록되어 있다''는 주장만 허용하고, 도출 증거는 ``정해진 규칙으로 계산했다''는 주장만 허용한다. 합성 가정과 합성 변이는 모델 체킹 절차를 시험하기 위한 입력이므로, 학습용 정상 궤적이나 실제 위험 사례로 바로 승격하지 않는다.

원천 근거가 있는 계약 통과(PASS)는 약한 유지 후보가 될 수 있지만, 합성 변이는 검사기 회귀 시험 전용이다. 기한이 독립 요구사항으로 정당화되지 않은 검토 후보(FAIL)는 학습용 위험 레이블이 아니다. 장면 간 연속성 증거가 없으면 출처 불충분으로 기록하고 장기 시퀀스 샘플로 묶지 않으며, 독립 클립/창 샘플로만 유지한다.
